\documentclass{article}%
\usepackage[T1]{fontenc}%
\usepackage[utf8]{inputenc}%
\usepackage{lmodern}%
\usepackage{textcomp}%
\usepackage{lastpage}%
\usepackage{geometry}%
\geometry{margin=2.2cm}%
%
\usepackage{amsmath}%
\usepackage{amssymb}%
\usepackage{booktabs}%
\usepackage{longtable}%
\usepackage{array}%
\usepackage{hyperref}%
\hypersetup{hidelinks}%
\title{IME 2027 {-} Física\\[4pt]\large Comparação com o gabarito Poliedro\thanks{Gerado com gpt-resolve: \url{https://github.com/lgabs/gpt-resolve}}}%
\date{28/10/2026}%
%
\begin{document}%
\normalsize%
\maketitle%
\section*{Resumo}\textbf{Nota média: 10,0/10} (30 de 30 itens com a resolução correta; cada item vale o mesmo dentro da questão).\par\medskip%
\begin{longtable}{@{}c c p{0.78\textwidth}@{}}\toprule Questão & Nota & Resumo \\ \midrule
1 & 10,0 & A resolução está correta e bate com o gabarito do Poliedro nos três itens: a) 250 m/s, b) 4000 K, c) 864 kJ, trabalho feito pelo gás. Refiz as contas em Python e conferi a transcrição com o PDF original. Os dados são coerentes entre si: c\_p/c\_v = 10/7 e massa molar de cerca de 28 g/mol. O único ponto discutível é a convenção de sinal do trabalho, que o enunciado não fixa. \\
2 & 10,0 & A resolução concorda com o gabarito Poliedro nos três itens. Em a), as reações são A = 200 kN para cima e B = 100 kN para baixo, sem componentes horizontais. Em b), os oito esforços normais têm os mesmos valores e a mesma natureza, e AE = 0. Em c), as duas respostas dão a mesma explicação: AE é a barra que impede o mecanismo de paralelogramo em ADEB. Conferi com python, resolvendo o sistema completo de 12 equações. A figura do PDF bate com a transcrição. Não encontrei erro nem na resolução nem no gabarito. \\
3 & 10,0 & A resolução está correta e coincide com o gabarito Poliedro nos três itens: a) 300 N para cima; b) 100 $^\circ$C; c) 37,5\% do dia, com Q = 5$\cdot$10$^{6}$ J e t = 10$^{4}$ s após o amanhecer. Refiz todas as contas e conferi a correção pela dilatação da área ($\approx$37,49\%), sem encontrar divergência. \\
4 & 10,0 & A resolução está correta e coincide com o gabarito Poliedro nos três itens. Item a: U\_AB = 900 V, i1 = 50 A, i2 = 0. Item b: razão 2,5, rompimento em $\approx$ 1,5 s. Item c: I\_cc = 500 A, t = 100 ms > 50 ms, o fusível não protege. Conferi as contas e medi o gráfico no PDF original, e os valores lidos ($\approx$ 1,1; 3,8 s; $\approx$ 3,35; 0,1 s) sustentam t $\approx$ 1,5 s. Os únicos pontos discutíveis são as leituras do gráfico e o prolongamento do patamar além de 5, e os dois resolvedores fazem a mesma coisa nos dois casos. \\
5 & 10,0 & A resolução está correta e coincide com o gabarito Poliedro nos três itens. Os resultados são T = 450 K, Q = 90 kW e, no item c, a conclusão de que a proposta B não é realista. B fornece no máximo 71,7 kW, suficientes para cerca de 215 residências, e não os 116,7 kW das 350. O aquecimento de água (33,3 kW, $\Delta$$\dot S$ $\approx$ +0,044 kW/K) é compatível com a 2ª lei. Recalculei todas as contas com frações exatas. A resolução traz ainda limites corretos que o gabarito não apresenta (Q$_{2}$ $\geq$ 20 kW e 85 kW reversível). \\
6 & 10,0 & A resolução bate com o gabarito Poliedro nos quatro itens: v = 68 m/s, R = 5$\surd$3 m, T $\approx$ 1,07$\cdot$10$^3$ N e V $\approx$ 55$\cdot$10$^{-}$$^3$ m$^3$. Conferi o enunciado no PDF original e refiz as contas em Python. A geometria do Doppler está correta: em A a fonte se aproxima em linha reta de O; em B ela se afasta, com cos$\theta$ = d/$\surd$(d$^2$+4R$^2$). A dinâmica também está: a componente horizontal da tração é a resultante centrípeta, e na vertical E = mg + T$\cdot$cos$\alpha$. Não encontrei erro em nenhum dos dois. \\
7 & 10,0 & A resolução está correta e coincide com o gabarito Poliedro nos três itens: a) 0,45 m; b) kx $\approx$ 6,44 N e razão $\approx$ -1,6 (29/18 em módulo); c) v = $\surd$2 m/s. A única diferença, no item b, é de convenção de sinal: a resolução dá o módulo e também o valor com sinal do gabarito. Conferi as contas com python. A hipótese de equilíbrio quase estático de P$_{2}$ no item b é a mesma nos dois e está declarada na resolução. \\
8 & 10,0 & A resolução e o gabarito Poliedro concordam nos quatro itens: (a) o anteparo está se afastando; (b) 0,53 m; (c) 0,254 m/s; (d) 1,792 m. Refiz as contas em Python com a tabela da prova e com as raízes exatas, que dão 0,529 m, 0,255 m/s e 1,795 m. As únicas diferenças vêm do arredondamento da tabela. Conferi o enunciado no PDF e ele bate com a transcrição. O método da resolução está correto: sin$\theta$\_m = m/4, a ordem 4 é rasante e não atinge o anteparo, e usa{-}se tg$\theta$ exata, sem aproximar para ângulos pequenos. A resolução está correta. \\
9 & 10,0 & A resolução da questão 9 está certa e bate com o gabarito do Poliedro nos dois itens. No item a, u = c($\surd$3h - L)/($\surd$3L - h). No item b, L/$\surd$3 < h < L. A geometria da reflexão (normal a 240$^\circ$, ângulo de incidência de 30$^\circ$) e as três condições (R$_{1}$ > 0, t$_{1}$ > 0, u < c) foram refeitas à mão e conferidas numericamente: o primeiro contato ocorre no t$_{1}$ previsto e o raio refletido sai em (0, -1). A resolução vai além do gabarito ao verificar que o contato é o primeiro e pela face externa. O único reparo é de redação, num trecho truncado dessa verificação. \\
10 & 10,0 & A resolução está correta e bate com o gabarito Poliedro nos dois itens. No item a), só a situação 2 serve, com justificativa para cada uma das quatro situações. No item b), o movimento horizontal é um MHS com k = 4$\mu$0 q v i/($\pi$ d$^2$), e p = v$\cdot$T/2 = ($\pi$ d/2)$\cdot$$\surd$($\pi$ m v/($\mu$0 q i)). Refiz os sentidos dos campos e das forças e conferi k, T e p numericamente, e tudo coincide. Não encontrei erro na resolução nem no gabarito. \\
\bottomrule\end{longtable}%
\section*{Questão 1}%
\href{https://poliedroresolve.sistemapoliedro.com.br/vestibulares/ime/2027/ime-2-fase-dia-2-27-10-2026-2-fase-dia-2/dissertativa-1-fisica-2-fase-dia-2-ime-2027}{Resolução do gabarito}\par\smallskip%
\begin{longtable}{@{}p{0.06\textwidth} p{0.33\textwidth} p{0.33\textwidth} p{0.1\textwidth} p{0.08\textwidth}@{}}\toprule Item & Resolução & Gabarito & Veredito & Certo \\ \midrule
a & v = 250 m/s & v0 = 250 m/s & igual & ambos \\
b & T = 4000 K & T\_falha = 4000 K & igual & ambos \\
c & W = 864 kJ (realizado pelo gás; Q = 2880 kJ) & W = 864 kJ (Q = 2880 kJ, $\Delta$U = 2016 kJ) & igual & ambos \\
\bottomrule\end{longtable}%
\begin{itemize}%
\item%
\textbf{a)} Refiz a conta pela equação da trajetória a 45$^\circ$: 1360 = 2000 - 10$\cdot$2000$^2$/v$^2$, o que dá v$^2$ = 62500 m$^2$/s$^2$ e v = 250 m/s. A resolução também confere que o foguete ainda sobe na vertical do alvo: v\_y = 45$\surd$2 $\approx$ 63,6 m/s, e o novo ápice fica em x = 3125 m, depois do alvo.%
\item%
\textbf{b)} Com T $\propto$ v$^2$, T = 2560$\cdot$(250/200)$^2$ = 2560$\cdot$25/16 = 4000 K.%
\item%
\textbf{c)} Q = m$\cdot$c\_p$\cdot$$\Delta$T = 2,0$\cdot$1,0$\cdot$1440 = 2880 kJ. Pela 1ª Lei, W = Q - $\Delta$U = 2880 - 2016 = 864 kJ. A resolução e o gabarito usam a mesma convenção: W é o trabalho feito pelo gás, positivo na expansão isobárica.%
\end{itemize}%
\paragraph{Pontos de dúvida}%
\begin{itemize}%
\item%
\textit{(statement)} "Trabalho realizado durante o aquecimento" não diz quem realiza o trabalho. A resolução e o gabarito tomam o trabalho feito pelo gás (+864 kJ). Se alguém tomar o trabalho feito sobre o gás, chega a -864 kJ. Vale ver se a correção oficial aceita as duas convenções, desde que declaradas.%
\item%
\textit{(solution)} Fica implícito que o foguete se comporta como um projétil livre desde que sai do lançador, sem empuxo depois do lançamento. É o modelo pedido pelo enunciado ('despreze a resistência do ar', velocidade de lançamento dada), mas a resolução poderia dizer isso de forma explícita.%
\end{itemize}%
\section*{Questão 2}%
\href{https://poliedroresolve.sistemapoliedro.com.br/vestibulares/ime/2027/ime-2-fase-dia-2-27-10-2026-2-fase-dia-2/dissertativa-2-fisica-2-fase-dia-2-ime-2027}{Resolução do gabarito}\par\smallskip%
\begin{longtable}{@{}p{0.06\textwidth} p{0.33\textwidth} p{0.33\textwidth} p{0.1\textwidth} p{0.08\textwidth}@{}}\toprule Item & Resolução & Gabarito & Veredito & Certo \\ \midrule
a & A: 200 kN, vertical, para cima; B: 100 kN, vertical, para baixo (apoio ancorado); A\_x = B\_x = 0 & A = 200 kN vertical para cima; B = 100 kN vertical para baixo (A\_x = B\_x = 0) & igual & ambos \\
b & CD = 75 kN (C), CF = 125 kN (T), DF = 200 kN (C), AD = 200 kN (C), AE = 0, DE = 75 kN (C), BE = 100 kN (T), EF = 125 kN (T) & N\_CD = 75 kN (C); N\_CF = 125 kN (T); N\_DF = 200 kN (C); N\_AD = 200 kN (C); N\_AE = 0; N\_DE = 75 kN (C); N\_BE = 100 kN (T); N\_EF = 125 kN (T) & igual & ambos \\
c & AE é barra de contraventamento: tem esforço nulo sob esta carga vertical, mas triangula o painel ADEB. Sem ela, ADEB seria um paralelogramo articulado (mecanismo com 1 grau de liberdade). Assim, AE dá rigidez e estabilidade lateral e resiste a cargas horizontais. & N\_AE = 0 só para este carregamento; AE é necessária para outras distribuições de carga e garante a rigidez (sem ela ADEB vira paralelogramo articulado e o guindaste balança lateralmente) & equivalente & ambos \\
\bottomrule\end{longtable}%
\begin{itemize}%
\item%
\textbf{a)} Refiz a conta. Em B chega só a barra vertical BE, então B\_x = 0, e pelo equilíbrio horizontal global A\_x = 0. Momentos em A: P$\cdot$3 + B\_y$\cdot$3 = 0, logo B\_y = -100 kN (para baixo). Daí A\_y = 200 kN (para cima). Resolvi também o sistema completo de 12 equações em python e ele confirma esses valores.%
\item%
\textbf{b)} Todos os valores e naturezas coincidem. Resolvi a treliça por matriz (método dos nós, 12 equações, posto 12) e obtive CD = -75, CF = +125, DF = -200, AD = -200, AE = 0, DE = -75, BE = +100 e EF = +125 (kN), com sinal positivo para tração. A resolução ainda confere as reações no nó A e no nó B.%
\item%
\textbf{c)} O conteúdo conceitual é o mesmo: esforço nulo neste carregamento, função de garantir rigidez e evitar o mecanismo de paralelogramo, e utilidade para outras cargas, como as horizontais. A resolução ainda justifica com a contagem m + r = 11 < 2j = 12 sem AE.%
\end{itemize}%
\paragraph{Pontos de dúvida}%
\begin{itemize}%
\item%
\textit{(statement)} Os dois apoios estão desenhados como articulações simples. A resposta, porém, exige que B puxe a estrutura para baixo (100 kN), ou seja, que o apoio esteja ancorado e resista a tração. A resolução declara isso, e o gabarito apenas dá o sentido para baixo. Vale confirmar se a banca espera essa observação explícita.%
\item%
\textit{(key)} A resolução do gabarito está resumida e tem atalhos, como 'A\_y + B\_y = 100' (em kN, sem unidade) e 'N\_BE = (4/5)N\_EF'. Esta última vale porque N\_AE = 0. As contas estão corretas, mas a fundamentação do item c) aparece só como 'ver item'.%
\end{itemize}%
\section*{Questão 3}%
\href{https://poliedroresolve.sistemapoliedro.com.br/vestibulares/ime/2027/ime-2-fase-dia-2-27-10-2026-2-fase-dia-2/dissertativa-3-fisica-2-fase-dia-2-ime-2027}{Resolução do gabarito}\par\smallskip%
\begin{longtable}{@{}p{0.06\textwidth} p{0.33\textwidth} p{0.33\textwidth} p{0.1\textwidth} p{0.08\textwidth}@{}}\toprule Item & Resolução & Gabarito & Veredito & Certo \\ \midrule
a & N = 300 N, vertical para cima (P = 600 N, T0 = 300 N) & R\_A = 300 N, vertical para cima (T = 300 N) & igual & ambos \\
b & $\Delta$$\theta$ = 100 $^\circ$C & $\Delta$T = 100 $^\circ$C & igual & ambos \\
c & Q = 5$\cdot$10$^{6}$ J; rompe 10$^{4}$ s após o amanhecer, às 37,5\% do dia ($\approx$37,49\% se considerar a dilatação da área) & Q = 5$\cdot$10$^{6}$ J; rompe às 37,5\% do dia (12,5\% após 25\%) & igual & ambos \\
\bottomrule\end{longtable}%
\begin{itemize}%
\item%
\textbf{a)} m = 2000$\cdot$(2$\cdot$0,25$\cdot$0,05) = 50 kg e P = 50$\cdot$12 = 600 N. O peso atua em L/2, então a reação no apoio e a tração ficam com P/2 = 300 N cada. A resolução e o gabarito coincidem.%
\item%
\textbf{b)} Tomando momentos no pino: T$\cdot$L = P$\cdot$L'/2, logo T = T0(1+$\alpha$$\Delta$$\theta$). A haste rompe com T = 1,001$\cdot$T0, o que dá $\alpha$$\Delta$$\theta$ = 10$^{-}$$^3$ e $\Delta$$\theta$ = 100 $^\circ$C. O raciocínio é o mesmo nos dois, e a conta foi conferida.%
\item%
\textbf{c)} A = 0,5 m$^2$. Na subida, I(t) = 0,2$\cdot$t (W/m$^2$, t em s), então E(t) = 0,05$\cdot$t$^2$. Igualando a 5$\cdot$10$^{6}$ J, t = 10$^{4}$ s = 12,5\% do dia, e o horário é 37,5\%. O gabarito trabalha com t em \% do dia e chega a 32000$\cdot$t$^2$ = 5$\cdot$10$^{6}$, ou seja, t = 12,5. Os dois métodos são equivalentes. A energia que chega até o meio{-}dia é 2$\cdot$10$^{7}$ J, maior que Q, então o rompimento ocorre mesmo na rampa de subida. Conferi a correção pela dilatação da área, integrando d$\theta$/(1+$\alpha$$\theta$)$^2$ = $\theta$/(1+$\alpha$$\theta$): dá cerca de 37,494\%, e o tempo muda cerca de 0,05\%. Isso confirma tanto o $\approx$37,49\% da resolução quanto o '0,05\%' citado pelo Poliedro.%
\end{itemize}%
\paragraph{Pontos de dúvida}%
\begin{itemize}%
\item%
\textit{(statement)} A dilatação da face superior (área (1+$\alpha$$\Delta$$\theta$)$^2$, até +0,2\%) é desprezada nas duas soluções. Se for considerada, o resultado vira $\approx$37,49\% em vez de 37,5\%. A diferença é desprezível, e tanto a resolução quanto o gabarito comentam isso.%
\item%
\textit{(statement)} Há hipóteses implícitas que ninguém discute: toda a energia absorvida vira aquecimento (sem perdas por emissão ou condução para os apoios), o calor específico é constante e a temperatura é uniforme na barra. Também se ignora o efeito da rotação do planeta sobre o peso: o raio R não é dado, então só dá para usar g = GM/R$^2$ = 12 m/s$^2$.%
\end{itemize}%
\section*{Questão 4}%
\href{https://poliedroresolve.sistemapoliedro.com.br/vestibulares/ime/2027/ime-2-fase-dia-2-27-10-2026-2-fase-dia-2/dissertativa-4-fisica-2-fase-dia-2-ime-2027}{Resolução do gabarito}\par\smallskip%
\begin{longtable}{@{}p{0.06\textwidth} p{0.33\textwidth} p{0.33\textwidth} p{0.1\textwidth} p{0.08\textwidth}@{}}\toprule Item & Resolução & Gabarito & Veredito & Certo \\ \midrule
a & U\_AB = 900 V; i1 = 50 A (Rec1); i2 = 0 (Rec2 no limiar, U\_AB igual à fcem dele); I = 50 A no gerador & U\_AB = 900 V; i1 = 50 A; i2 = 0 (tensão A--B igual à fcem de Rec2) & igual & ambos \\
b & I/I\_FUS = 2,5, no trecho inclinado da curva; o fusível se rompe cerca de 1,5 s após o fechamento de S & I/I\_fus = 2,5 $\Rightarrow$ o fusível rompe em t $\approx$ 1,5 s (lido na curva) & igual & ambos \\
c & I\_cc = 1000/2 = 500 A (25$\cdot$I\_FUS); a curva dá o patamar mínimo de 0,1 s = 100 ms > 50 ms; o fusível NÃO protege a fonte & I\_curto = 500 A, I/I\_fus = 25 $\Rightarrow$ t = 100 ms > 50 ms $\Rightarrow$ o fusível NÃO protege a fonte & igual & ambos \\
\bottomrule\end{longtable}%
\begin{itemize}%
\item%
\textbf{a)} Refiz pelo teorema de Millman: U = (500+850+450)/(0,5+1+0,5) = 900 V. Daí i1 = (900-850)/1 = 50 A, i2 = (900-900)/2 = 0 e I = (1000-900)/2 = 50 A = i1+i2. O balanço de potências da resolução (50 kW = 5 kW + 45 kW) também confere.%
\item%
\textbf{b)} I = 50 A e 50/20 = 2,5. Medi a figura no PDF renderizado: o trecho vertical fica em x $\approx$ 1,1 e termina em t $\approx$ 3,8 s, e a reta chega a 0,1 s em x $\approx$ 3,35. Interpolando de (1,1; 3,8) a (3,4; 0,1), dá t $\approx$ 1,55 s. Medindo direto os pixels em x = 2,5, dá t $\approx$ 1,51 s. Se o fim da reta for tomado em 3,3, como sugere o gabarito, dá $\approx$ 1,45 s. Em todos os casos o resultado é $\approx$ 1,5 s.%
\item%
\textbf{c)} Com A e B em curto, U\_AB = 0 e a corrente no fusível é $\varepsilon$/r = 500 A, razão 25. Acima de $\approx$ 3,4 a curva é horizontal em 10\^{}{-}1 s = 100 ms, e isso é maior que 50 ms, então o fusível não protege. A resolução também observa, corretamente, que uma eventual corrente devolvida pelos receptores se fecharia pelo curto e não passaria por F.%
\end{itemize}%
\paragraph{Pontos de dúvida}%
\begin{itemize}%
\item%
\textit{(statement)} Para a razão 25, o gráfico da Figura 2 só vai até I/I\_FUS $\approx$ 5. A resolução e o gabarito prolongam o patamar horizontal de 0,1 s além desse limite. É a leitura natural, mas não está escrita no enunciado.%
\item%
\textit{(solution)} Os pontos (1,1; 3,8 s) e (3,4; 0,1 s) foram lidos no gráfico. O gabarito fala em $\approx$ 3,3 para o fim da reta. Na figura renderizada medi $\approx$ 3,35. A variação no tempo de rompimento do item b é de $\pm$0,1 s, e o resultado $\approx$ 1,5 s continua valendo.%
\item%
\textit{(statement)} Os dois resolvedores supõem que as fcem dos receptores são constantes desde t = 0, o que dá corrente constante de 50 A até o rompimento. Se os receptores fossem motores partindo do repouso, a fcem inicial seria nula e a corrente de partida seria maior. O enunciado fornece as fcem, então a hipótese adotada é a esperada, e a resolução a declara explicitamente.%
\item%
\textit{(statement)} O enunciado diz que a Figura 2 está em escala linear, mas o eixo vertical tem o rótulo 10\^{}{-}1 na posição da origem. Isso é ambíguo, mas não muda a leitura do patamar (0,1 s = 100 ms).%
\end{itemize}%
\section*{Questão 5}%
\href{https://poliedroresolve.sistemapoliedro.com.br/vestibulares/ime/2027/ime-2-fase-dia-2-27-10-2026-2-fase-dia-2/dissertativa-5-fisica-2-fase-dia-2-ime-2027}{Resolução do gabarito}\par\smallskip%
\begin{longtable}{@{}p{0.06\textwidth} p{0.33\textwidth} p{0.33\textwidth} p{0.1\textwidth} p{0.08\textwidth}@{}}\toprule Item & Resolução & Gabarito & Veredito & Certo \\ \midrule
a & $\eta$\_A = 0,40, $\eta$\_Carnot = 0,50 e T = 450 K & $\eta$1 = 40\%; T = 450 K & igual & ambos \\
b & Q = 150 - 60 = 90 kW & Q rejeitado = 90 kW & igual & ambos \\
c & Não é realista. P necessária = 116,7 kW e Q\_água = 33,3 kW. Pela 1ª lei, $\dot W$\_B = 150 - 45 - 33,3 $\approx$ 71,7 kW, o que atende só 215 residências. O aquecimento é viável pela 2ª lei: $\dot S$\_ger = 2/45 $\approx$ 0,044 kW/K > 0, e Q$_{2}$ precisa ser de pelo menos 20 kW. Mesmo reversível, a máquina daria no máximo 85 kW. & P = 116,7 kW; Q\_água = 33,3 kW; $\dot W$ = 71,7 kW < 116,7 kW, logo não supre as 350 residências (máx. 215); $\Delta$$\dot S$\_universo = +0,04 kW/K, logo não viola a 2ª lei & igual & ambos \\
\bottomrule\end{longtable}%
\begin{itemize}%
\item%
\textbf{a)} Como 80 HP = 60 kW, temos $\eta$ = 60/150 = 0,4. Então 0,4 = 0,8(1 - T/900), o que dá T = 450 K. Conferi com frações exatas.%
\item%
\textbf{b)} Sai da 1ª lei aplicada ao ciclo. A resolução também verifica que $\dot S$\_ger = 1/30 kW/K > 0, o que é coerente.%
\item%
\textbf{c)} Todos os números coincidem: 116,7 kW, 33,3 kW, 71,7 kW, 215 residências e $\Delta$$\dot S$ = 2/45 $\approx$ 0,044 kW/K (o gabarito arredonda para 0,04). A conclusão também é a mesma: a energia elétrica para 350 casas não fecha pela 1ª lei, e o aquecimento de água é compatível com a 2ª lei. A resolução acrescenta dois limites que o gabarito não traz: Q$_{2}$ $\geq$ 20 kW e $\dot W$\_máx reversível = 85 kW, que corresponde a 255 casas. Estão corretos e não contradizem o gabarito.%
\end{itemize}%
\paragraph{Pontos de dúvida}%
\begin{itemize}%
\item%
\textit{(solution)} A frase "superaria até o limite reversível da 2ª lei (85 kW)" mistura dois cenários. Os 85 kW valem para Q$_{2}$ = 20 kW, que não basta para aquecer 50 L/min. Com Q$_{2}$ = 33,3 kW, o teto é 71,7 kW pela 1ª lei. A observação é correta como limite superior, mas o argumento decisivo é a 1ª lei.%
\item%
\textit{(key)} O gabarito arredonda $\Delta$$\dot S$ para +0,04 kW/K. O valor exato é 2/45 $\approx$ 0,044 kW/K. É só arredondamento, não é erro.%
\item%
\textit{(statement)} O enunciado não diz se a máquina B tem a mesma eficiência relativa (80\% de Carnot) que a A. As duas soluções tomam como dados só Q\_H = 150 kW e o calor de 45 kW rejeitado a 450 K, e supõem que toda a potência útil vira eletricidade. Essa é a leitura natural.%
\end{itemize}%
\section*{Questão 6}%
\href{https://poliedroresolve.sistemapoliedro.com.br/vestibulares/ime/2027/ime-2-fase-dia-2-27-10-2026-2-fase-dia-2/dissertativa-6-fisica-2-fase-dia-2-ime-2027}{Resolução do gabarito}\par\smallskip%
\begin{longtable}{@{}p{0.06\textwidth} p{0.33\textwidth} p{0.33\textwidth} p{0.1\textwidth} p{0.08\textwidth}@{}}\toprule Item & Resolução & Gabarito & Veredito & Certo \\ \midrule
a & v = 68 m/s & v = 68 m/s & igual & ambos \\
b & R = 5$\surd$3 $\approx$ 8,66 m (cos$\theta$ = 1/2) & $\theta$ = 60$^\circ$ (cos$\theta$ = 1/2); R = 5$\surd$3 m & igual & ambos \\
c & T $\approx$ 1067,8 N $\approx$ 1,07$\cdot$10$^3$ N ($\alpha$ = 60$^\circ$) & $\alpha$ = 60$^\circ$; T $\cong$ 1,07$\cdot$10$^3$ N & igual & ambos \\
d & V $\approx$ 5,51$\cdot$10$^{-}$$^2$ m$^3$ ($\approx$ 55 L), com E $\approx$ 551,2 N & E $\approx$ 551 N; V $\cong$ 55$\cdot$10$^{-}$$^3$ m$^3$ & igual & ambos \\
\bottomrule\end{longtable}%
\begin{itemize}%
\item%
\textbf{a)} Em A a velocidade aponta direto para O: 5/4 = 340/(340-v), logo v = 68 m/s. A resolução e o gabarito batem.%
\item%
\textbf{b)} Em B a fonte se afasta: 10/11 = 340/(340 + v$\cdot$cos$\theta$), logo v$\cdot$cos$\theta$ = 34 e cos$\theta$ = 1/2. A resolução usa cos$\theta$ = d/$\surd$(d$^2$+4R$^2$). O gabarito usa tg$\theta$ = 2R/d. As duas relações são equivalentes e dão R = ($\surd$3/2)$\cdot$d = 5$\surd$3 m.%
\item%
\textbf{c)} Temos sen$\alpha$ = R/L = $\surd$3/2 e T$\cdot$sen$\alpha$ = mv$^2$/R. Daí T = (8/5)$\cdot$(34$^2$/$\surd$3) = 1,6$\cdot$667,4 $\approx$ 1067,8 N. A conta exata em Python dá 1067,87 N; a diferença vem só do arredondamento do dado 667,4.%
\item%
\textbf{d)} Na vertical, E = mg + T$\cdot$cos$\alpha$ = 17,32 + 533,9 $\approx$ 551,2 N. Então V = E/($\rho$g) $\approx$ 0,0551 m$^3$ = 55$\cdot$10$^{-}$$^3$ m$^3$. Os dois chegam ao mesmo valor, só escrito de outra forma.%
\end{itemize}%
\section*{Questão 7}%
\href{https://poliedroresolve.sistemapoliedro.com.br/vestibulares/ime/2027/ime-2-fase-dia-2-27-10-2026-2-fase-dia-2/dissertativa-7-fisica-2-fase-dia-2-ime-2027}{Resolução do gabarito}\par\smallskip%
\begin{longtable}{@{}p{0.06\textwidth} p{0.33\textwidth} p{0.33\textwidth} p{0.1\textwidth} p{0.08\textwidth}@{}}\toprule Item & Resolução & Gabarito & Veredito & Certo \\ \midrule
a & 0,45 m & d = 0,45 m & igual & ambos \\
b & E\_el/|$\Delta$U\_g| = 29/18 $\approx$ 1,6 (com sinal, $\approx$ -1,6); kx = 58/9 $\approx$ 6,44 N & kx $\cong$ 6,44 N; U\_el/$\Delta$U\_g $\cong$ -1,6 & equivalente & ambos \\
c & v = $\surd$2 $\approx$ 1,41 m/s (P$_{2}$ atinge o vértice C, velocidade dirigida de B para C) & v2 = $\surd$2 m/s & igual & ambos \\
\bottomrule\end{longtable}%
\begin{itemize}%
\item%
\textbf{a)} Com P$_{2}$ fixa, a conservação da energia de P$_{1}$ dá 2(1-r) = 0,9(1/r - 1). Descartando a raiz r = 1, fica r = 0,9/2 = 0,45 m. Os dois chegam a isso pelo mesmo caminho, e conferi a conta.%
\item%
\textbf{b)} Os dois usam o mesmo modelo: kx = mg$\cdot$sen30$^\circ$ + K\_e$\cdot$q$_{1}$q$_{2}$/r$^2$ = 2 + 4,44 = 6,44 N (58/9 N), com E\_el = kx$\cdot$x/2 e $\Delta$U\_g = -mg$\cdot$sen30$^\circ$$\cdot$x. Daí a razão é -kx/4 $\approx$ -1,61. A resolução dá o módulo, 29/18 $\approx$ 1,6, e também escreve o valor com sinal, -1,6, igual ao gabarito. A diferença é só de convenção de sinal.%
\item%
\textbf{c)} Depois do giro de 30$^\circ$ no sentido anti{-}horário, BC fica horizontal, como mostra a Figura 3. Como a quantidade de movimento se conserva e as massas são iguais, as duas partículas andam distâncias iguais em sentidos opostos. P$_{2}$ anda 0,15 m até C, enquanto P$_{1}$ ainda está sobre BC, a 0,40 m de B, e a separação vai a 0,75 m. A energia liberada é 0,9(1/0,45 - 1/0,75) = 0,8 J, e 0,8 = 0,4v$^2$ dá v = $\surd$2 m/s. Conferi com python. A resolução ainda mostra que C é o primeiro vértice atingido, coisa que o gabarito só deixa implícita.%
\end{itemize}%
\paragraph{Pontos de dúvida}%
\begin{itemize}%
\item%
\textit{(statement)} Item b: a mola começa relaxada, mas sobre P$_{2}$ já atuam o peso tangencial (2 N) e a repulsão inicial (0,9 N), então P$_{2}$ não está em equilíbrio no começo. A resolução e o gabarito tratam P$_{2}$ como em equilíbrio quase estático na aproximação máxima (kx = 6,44 N) e desprezam a oscilação de P$_{2}$. Com essa oscilação, a compressão máxima chegaria a cerca do dobro da estática, e a razão mudaria. O termo 'estime' justifica a hipótese, mas ela não está escrita no enunciado.%
\item%
\textit{(statement)} Item b, sinal: 'razão entre a energia potencial elástica e a variação da energia potencial gravitacional' admite duas leituras, com sinal (-1,6, como no gabarito) ou em módulo (1,6). A resolução dá as duas formas.%
\item%
\textit{(statement)} Item c: o enunciado não diz se o anteparo continua em C depois que a mola é retirada. Isso não muda a velocidade de P$_{2}$ ao chegar a C, que é $\surd$2 m/s nos dois casos.%
\end{itemize}%
\section*{Questão 8}%
\href{https://poliedroresolve.sistemapoliedro.com.br/vestibulares/ime/2027/ime-2-fase-dia-2-27-10-2026-2-fase-dia-2/dissertativa-8-fisica-2-fase-dia-2-ime-2027}{Resolução do gabarito}\par\smallskip%
\begin{longtable}{@{}p{0.06\textwidth} p{0.33\textwidth} p{0.33\textwidth} p{0.1\textwidth} p{0.08\textwidth}@{}}\toprule Item & Resolução & Gabarito & Veredito & Certo \\ \midrule
a & O anteparo está se afastando da rede: y\_m = D$\cdot$tg$\theta$\_m cresce com D, e as ordens mais altas saem pelas bordas. & O anteparo está se afastando da rede & igual & ambos \\
b & D(t1) = 0,2$\surd$7 $\approx$ 0,53 m & D1 = 0,53 m & igual & ambos \\
c & v $\approx$ 0,254 m/s ($\approx$ 0,25 m/s) & v = 0,254 m/s & igual & ambos \\
d & $\Delta$s = 0,6$\surd$15 - 0,2$\surd$7 $\approx$ 1,79 m (1,792 m) & $\Delta$s = 1,792 m & igual & ambos \\
\bottomrule\end{longtable}%
\begin{itemize}%
\item%
\textbf{a)} Mesma conclusão e mesma justificativa. Como $\theta$\_m é fixo, a altura y\_m = D$\cdot$tg$\theta$\_m só cresce se D aumenta, e o número de máximos visíveis só cai (7$\to$5$\to$3) nesse caso.%
\item%
\textbf{b)} sin$\theta$3 = 3/4 dá tg$\theta$3 = 3/$\surd$7, e D = 0,6$\cdot$$\surd$7/3 = 0,2$\surd$7. Pela tabela, 0,2$\cdot$2,65 = 0,53 m; o valor exato é 0,529 m. O gabarito chama essa distância de D1, nome diferente do usado na resolução, mas o valor é o mesmo.%
\item%
\textbf{c)} D2 = 0,6$\surd$3, que dá 1,038 m pela tabela. Então v = (1,038 - 0,530)/2 = 0,254 m/s. Com as raízes exatas dá 0,255 m/s; a diferença vem só do arredondamento da tabela.%
\item%
\textbf{d)} O par $\pm$1 sai das bordas em D = 0,6$\surd$15, que dá 2,322 m pela tabela. Então $\Delta$s = 2,322 - 0,530 = 1,792 m, igual ao gabarito; o valor exato é 1,795 m. A resolução escreve 1,792 m na conta e arredonda para 1,79 m na resposta, o que é equivalente.%
\end{itemize}%
\section*{Questão 9}%
\href{https://poliedroresolve.sistemapoliedro.com.br/vestibulares/ime/2027/ime-2-fase-dia-2-27-10-2026-2-fase-dia-2/dissertativa-9-fisica-2-fase-dia-2-ime-2027}{Resolução do gabarito}\par\smallskip%
\begin{longtable}{@{}p{0.06\textwidth} p{0.33\textwidth} p{0.33\textwidth} p{0.1\textwidth} p{0.08\textwidth}@{}}\toprule Item & Resolução & Gabarito & Veredito & Certo \\ \midrule
a & u = c($\surd$3h - L)/($\surd$3L - h), com t$_{1}$ = ($\surd$3L - h)/c e R$_{1}$ = $\surd$3h - L & u = c($\surd$3$\cdot$h - L)/($\surd$3$\cdot$L - h) & igual & ambos \\
b & L/$\surd$3 < h < L, isto é, tan30$^\circ$ < h/L < tan45$^\circ$ & L/$\surd$3 < h < L & igual & ambos \\
\bottomrule\end{longtable}%
\begin{itemize}%
\item%
\textbf{a)} A expressão é a mesma. Refiz a conta: o raio incidente está a 30$^\circ$ e o refletido a 270$^\circ$, então a normal externa no ponto de impacto aponta a 240$^\circ$, com ângulo de incidência de 30$^\circ$. Daí P = (L - R$_{1}$/2, h - $\surd$3R$_{1}$/2). Impondo que P esteja sobre a reta x = $\surd$3y, sai R$_{1}$ = $\surd$3h - L e OP = 2y\_P = $\surd$3L - h. Esse é o mesmo sistema do gabarito (ct$\cdot$$\surd$3/2 + ut/2 = L; ct/2 + ut$\cdot$$\surd$3/2 = h). Simulei numericamente para (L,h) = (1; 0,8), (1; 0,6) e (2; 1,9). Em todos os casos o primeiro contato ocorre no t$_{1}$ previsto e o raio refletido sai em (0, -1). A resolução também confere que o contato é o primeiro e ocorre pela face externa (d$\cdot$n < 0 e f'(t$_{1}$) < 0).%
\item%
\textbf{b)} As três exigências são as mesmas nas duas: R$_{1}$ > 0 dá h > L/$\surd$3; t$_{1}$ > 0 dá h < $\surd$3L; u < c dá ($\surd$3+1)h < ($\surd$3+1)L, ou seja, h < L. Como h < L já implica h < $\surd$3L, a condição final coincide. A resolução ainda dá a leitura geométrica: visto da origem, o centro fica entre as retas de 30$^\circ$ e 45$^\circ$.%
\end{itemize}%
\paragraph{Pontos de dúvida}%
\begin{itemize}%
\item%
\textit{(solution)} Na verificação do primeiro contato, a frase "Então f$^2$-\ldots{} leva a" está truncada e pouco clara. O argumento continua válido, mas está mal redigido: g($\tau$) = |c$\tau$d - C|$^2$ - u$^2$$\tau$$^2$ é uma parábola com concavidade para cima, e g' < 0 em t$_{1}$ mostra que t$_{1}$ é a menor raiz. Vale reescrever esse trecho.%
\item%
\textit{(statement)} O enunciado não diz se o caso{-}limite h = L/$\surd$3 (R$_{1}$ = 0, u = 0, o centro sobre a trajetória) é aceito. Tanto a resolução quanto o gabarito usam desigualdades estritas, o que é razoável, já que nesse caso não há reflexão numa superfície de raio positivo.%
\end{itemize}%
\section*{Questão 10}%
\href{https://poliedroresolve.sistemapoliedro.com.br/vestibulares/ime/2027/ime-2-fase-dia-2-27-10-2026-2-fase-dia-2/dissertativa-10-fisica-2-fase-dia-2-ime-2027}{Resolução do gabarito}\par\smallskip%
\begin{longtable}{@{}p{0.06\textwidth} p{0.33\textwidth} p{0.33\textwidth} p{0.1\textwidth} p{0.08\textwidth}@{}}\toprule Item & Resolução & Gabarito & Veredito & Certo \\ \midrule
a & Só a situação 2. Nas situações 1 e 4, B $\neq$ 0 no eixo, então há força lateral já em x = 0 e não existe equilíbrio. Na situação 3, o equilíbrio é instável. Na situação 2, F\_x $\approx$ -(4$\mu$0 q v i/($\pi$ d$^2$))$\cdot$x, que é restauradora. & Somente a situação 2 & igual & ambos \\
b & $\omega$$^2$ = 4$\mu$0 q v i/($\pi$ m d$^2$), T = $\pi$ d $\surd$($\pi$ m/($\mu$0 q v i)), p = v$\cdot$T/2 = ($\pi$ d/2)$\cdot$$\surd$($\pi$ m v/($\mu$0 q i)) & T = $\pi$d$\cdot$$\surd$(m$\pi$/($\mu$0 i q v)); p = ($\pi$d/2)$\cdot$$\surd$($\pi$mv/($\mu$0 i q)) & igual & ambos \\
\bottomrule\end{longtable}%
\begin{itemize}%
\item%
\textbf{a)} Refiz os sentidos com os fios em x = $\mp$d/2 e as correntes em +$\hat y$. O fio esquerdo gera -$\hat z$ à sua direita e o direito gera +$\hat z$ à sua esquerda, então B\_z = ($\mu$0 i/2$\pi$){[}1/(d/2-x) - 1/(d/2+x){]} $\approx$ 4$\mu$0 i x/($\pi$ d$^2$). Com v = -v $\hat y$, a força dá F\_x = -q v B\_z, que é restauradora (situação 2). Com v para cima, o sinal se inverte e o equilíbrio fica instável (situação 3). Com correntes opostas, os dois campos se somam entre os fios, |B| $\geq$ 2$\mu$0 i/($\pi$ d) no eixo, e não há equilíbrio (situações 1 e 4). A resolução justifica cada uma das quatro situações, como o item pede, e chega à mesma conclusão do gabarito.%
\item%
\textbf{b)} Conferi numericamente com um próton, v = 1e5 m/s, i = 10 A e d = 1 cm. A constante k obtida por derivada numérica de q v B\_z bate com 4$\mu$0 q v i/($\pi$ d$^2$) (2,56e{-}15 N/m). Também batem T $\approx$ 5,07 $\mu$s e p $\approx$ 0,254 m. A figura mostra p entre dois cruzamentos consecutivos com o eixo y, o que corresponde a meio período, e foi isso que a resolução e o gabarito usaram. A resolução também argumenta corretamente que v\_y $\approx$ -v e que F\_y = -q v\_x B\_z são correções de 2ª ordem.%
\end{itemize}%
\paragraph{Pontos de dúvida}%
\begin{itemize}%
\item%
\textit{(statement)} O enunciado não diz se o peso deve ser desprezado. A resolução assume isso explicitamente, e o gabarito também, implicitamente. Para uma partícula carregada com v constante é a hipótese natural, mas não está escrita.%
\end{itemize}%
\end{document}